% =====================================================================
%  4 Results I: behaviour (RQ1, RQ2)
%  Sources: reports/v3_identity_vs_baseline.md 5 (standard metrics),
%  reports/v2_behavioural_analysis.md (sections 1-9),
%  reports/v3_runB_results.md 1 (place controls),
%  S/results/v2/analysis/behaviour_*.csv, v3_region_vs_place.csv.
%  Numbers in the tables are copied from those reports: re-check against
%  the CSVs before submission.
% =====================================================================
\section{Results I: How Demographic Framing Changes Judgments}
\label{sec:behaviour}

% ---------------------------------------------------------------------
\subsection{Unframed baseline}
\label{sec:behaviour:baseline}
% Both models rank texts reasonably (AUC > 0.8) but over-flag:
% recall 0.99, precision about 0.26. Over-flagging matters later: any
% method that lowers P(yes) "improves" distance to raters.
\fillin{Baseline performance on all texts; over-flagging.}

\begin{table}[t]
  \centering
  \caption{Unframed prompt on all \num{4519} \dthree{} texts. Majority label: at least half
    of the raters call the text offensive (21.7\% of texts).}
  \label{tab:standard-metrics}
  \src{v3\_identity\_vs\_baseline.md §5}
  \begin{tabular}{lcc}
    \toprule
                              & \llamashort & \qwenshort \\
    \midrule
    texts flagged             & 78.7\% & 85.0\% \\
    accuracy                  & 0.42 & 0.37 \\
    precision                 & 0.27 & 0.25 \\
    recall                    & 0.99 & 0.99 \\
    F1 (offensive)            & 0.43 & 0.40 \\
    macro-F1                  & 0.42 & 0.36 \\
    ROC-AUC                   & 0.83 & 0.82 \\
    Brier vs raters' share    & 0.26 & 0.35 \\
    Spearman vs raters' share & 0.66 & 0.67 \\
    \bottomrule
  \end{tabular}
\end{table}

% ---------------------------------------------------------------------
\subsection{Demographic framings against placebos}
\label{sec:behaviour:placebo}
% - Variance between phrases is as large as between items.
% - Placebos: negative length slope in every variant; Qwen placebos push
%   towards "offensive" (+2.3 to +3.8), Llama placebos near 0; tea outlier.
% - Region beats the length-adjusted placebo line in all 10 variants.
%   Gender, religion, bare identity: not distinguishable (low power with
%   2 phrases: say "not distinguishable", not "no effect").
% - Exception: Llama context variants, gender +1.07 (p = 0.007).
% - Single regions: Sub-Saharan Africa strongest in both models.
\fillin{Placebo behaviour, then the headline comparison.}

\begin{table}[t]
  \centering
  \caption{Demographic framings minus content-free placebos, adjusted for phrase length
    (\variant{prefix}, log-odds). $p$-values treat the phrase as the unit.}
  \label{tab:identity-vs-placebo}
  \src{v2\_behavioural\_analysis.md §2}
  \begin{tabular}{lrrrr}
    \toprule
    & \multicolumn{2}{c}{\llamashort} & \multicolumn{2}{c}{\qwenshort} \\
    \cmidrule(lr){2-3}\cmidrule(lr){4-5}
    axis (phrases)           & diff & $p$ & diff & $p$ \\
    \midrule
    all demographics (14)    & $+1.21$ & $<0.001$ & $+2.38$ & $<0.001$ \\
    region (8)               & $\bm{+1.62}$ & $\bm{<0.001}$ & $\bm{+2.78}$ & $\bm{0.001}$ \\
    gender (2)               & $+0.22$ & 0.59 & $+0.27$ & 0.79 \\
    religion (2)             & $-0.10$ & 0.80 & $+1.17$ & 0.33 \\
    bare identity (2)        & $+0.70$ & 0.13 & $+0.96$ & 0.34 \\
    \bottomrule
  \end{tabular}
\end{table}

% Optional second table / figure: per-region effects (v2 report §2), or a
% figure of the shift per condition with the placebo length line.
\begin{figure}[t]
  \centering
  \placeholderfig{Shift per condition vs phrase length, placebo regression line,
    demographics coloured by axis; one panel per model (\variant{prefix}).}
  \caption{\fillin{Caption.}}
  \label{fig:length-null}
\end{figure}

% ---------------------------------------------------------------------
\subsection{Prompt format}
\label{sec:behaviour:format}
% - Paraphrases agree (rho 0.83 to 0.99).
% - Suffix pulls everything down; in Qwen it reverses the identity effect
%   (-2.2 log-odds).
% - Target measures topic relevance (placebos collapse), report it
%   separately from perspective formats.
% - Context makes Llama use gender.
% - Flip rates: Llama 29% vs 21%, Qwen 38% vs 34% (prefix).
\fillin{How the format changes size and even sign of the effect.}

\begin{table}[t]
  \centering
  \caption{Mean \pyes{} for demographic and placebo framings per format
    (unframed mean: \llamashort{} 0.47, \qwenshort{} 0.53).}
  \label{tab:formats}
  \src{v2\_behavioural\_analysis.md §5}
  \begin{tabular}{lcccc}
    \toprule
    & \multicolumn{2}{c}{\llamashort} & \multicolumn{2}{c}{\qwenshort} \\
    \cmidrule(lr){2-3}\cmidrule(lr){4-5}
    format & demographic & placebo & demographic & placebo \\
    \midrule
    \variant{prefix}        & 0.66 & 0.47 & 0.87 & 0.70 \\
    \variant{infix}         & 0.64 & 0.39 & 0.76 & 0.58 \\
    \variant{suffix}        & 0.24 & 0.17 & 0.22 & 0.53 \\
    \variant{embedded}      & 0.50 & 0.27 & 0.65 & 0.50 \\
    \variant{target}        & 0.50 & 0.16 & 0.51 & 0.15 \\
    \variant{context}       & 0.74 & 0.54 & 0.75 & 0.65 \\
    \variant{context\_long} & 0.67 & 0.49 & 0.84 & 0.77 \\
    \bottomrule
  \end{tabular}
\end{table}

% ---------------------------------------------------------------------
\subsection{Identity or place name?}
\label{sec:behaviour:place}
% Region minus its two place controls ("researching / mapping X").
% Llama: about half is identity (prefix +0.73, infix +1.01).
% Qwen: perspective formats are the place name (prefix +0.02);
% identity part only after the text (infix +1.17, embedded +1.16).
\fillin{Place-name controls and what they change about the region result.}

\begin{table}[t]
  \centering
  \caption{Region framing minus its two place-name controls (log-odds, mean over 4 regions;
    $p$ with the phrase as the unit).}
  \label{tab:place}
  \src{v3\_runB\_results.md §1}
  \begin{tabular}{lrrrr}
    \toprule
    & \multicolumn{2}{c}{\llamashort} & \multicolumn{2}{c}{\qwenshort} \\
    \cmidrule(lr){2-3}\cmidrule(lr){4-5}
    format & diff & $p$ & diff & $p$ \\
    \midrule
    \variant{prefix}   & $+0.73$ & 0.013 & $+0.02$ & 0.92 \\
    \variant{infix}    & $+1.01$ & 0.003 & $+1.17$ & 0.007 \\
    \variant{suffix}   & $-0.05$ & 0.67  & $-0.02$ & 0.90 \\
    \variant{embedded} & $+0.45$ & 0.07  & $+1.16$ & 0.016 \\
    \variant{target}   & $+0.75$ & 0.11  & $+5.09$ & 0.02 \\
    \variant{context}  & $+0.76$ & 0.06  & $-0.38$ & 0.03 \\
    \bottomrule
  \end{tabular}
\end{table}

% ---------------------------------------------------------------------
\subsection{Do the shifts match the real groups?}
\label{sec:behaviour:alignment}
% RQ2.
% - Absolute distance (prefix): Llama 0.35 -> 0.51 framed (0.34 matched
%   placebo); Qwen 0.40 -> 0.68 (0.50).
% - Relative: 0 of 240 tests significant; mean Spearman 0.02 / 0.00.
% - Region order: Llama matches the stable average offsets (rho 0.81,
%   n = 8) with twice the spread; Qwen 0.05 (v3_identity_vs_baseline.md 4).
%   This replaces the v2 value 0.36 that used noisy per-text gaps.
\fillin{Absolute and relative alignment with real group ratings.}
\src{v2 §8; v3\_identity §4}
